\subsection{Small batches change the competitive region}
\label{sec:7-smallbatch}

T5 compares scheduled JHQ, CAGRA fp32/int8 and IVF-RaBitQ at batches
8/32 on the four E1 datasets. Five synchronous host-to-host rounds use
1,000/992 queries, respectively. Its 1,808 points include JHQ's one-block
control; IVF-PQ and brute force are outside this campaign.
\Cref{fig:7-smallbatch,tab:7-smallbatch} compare the resulting regions.

\suppressfloats[t]
\begin{figure}[htbp]
\centering
\includegraphics[width=\linewidth]{t4-t5/smallbatch.pdf}
\caption[Small-batch competitive regions]{T5 JHQ / fastest-competitor
QPS under interpolation with qualifying measured endpoints. Above one
favours JHQ. \Cref{tab:7-smallbatch} gives strict thresholds.}
\label{fig:7-smallbatch}
\end{figure}

\begin{table}[htbp]
\centering
\caption[Small-batch matched-recall ratios]{T5 QPS ratios:
interpolation with qualifying endpoints / strict measured threshold.
The fastest competitor is selected independently for each estimator;
higher-recall qualifying endpoints remain eligible without extrapolation.}
\label{tab:7-smallbatch}
\footnotesize\setlength{\tabcolsep}{4pt}
\begin{tabular}{@{}lrrrr@{}}
\toprule
Dataset & $b$ & $R=0.90$ & $R=0.95$ & $R=0.97$ \\
\midrule
Vogue & 8 & 1.155 / 1.091 & 1.073 / 0.980 & 0.893 / 0.854 \\
Vogue & 32 & 1.069 / 0.997 & 0.956 / 0.895 & 0.865 / 0.839 \\
arXiv & 8 & 1.479 / 1.454 & 0.859 / 0.851 & 0.697 / 0.698 \\
arXiv & 32 & 1.166 / 1.148 & 0.667 / 0.661 & 0.622 / 0.582 \\
OpenAI-1536 & 8 & 1.316 / 1.131 & 0.851 / 0.825 & 0.641 / 0.620 \\
OpenAI-1536 & 32 & 1.071 / 0.915 & 0.770 / 0.752 & 0.937 / 0.914 \\
OpenAI-3072 & 8 & 1.032 / 0.953 & 0.798 / 0.798 & 0.721 / 0.702 \\
OpenAI-3072 & 32 & 0.866 / 0.734 & 0.608 / 0.771 & 1.008 / 1.351 \\
\bottomrule
\end{tabular}

\end{table}

At recall 0.97 and batch 8, CAGRA fp32 leads on all four datasets under
both estimators: interpolated JHQ ratios are 0.641--0.893. Batch-32 ratios are 0.622--1.008: OpenAI-3072 is near parity under interpolation, JHQ-favourable under the measured threshold. Vogue's batch-8/0.95 lead also depends on interpolation. Borderline rankings depend on grid resolution.

Qualifying endpoints above the target remain eligible. At recall 0.90, batch-8/32 ratios are 1.032--1.479/0.866--1.166, with three/one strict-threshold JHQ wins. Dropping endpoints overstates leads. Scheduling gains $1.7$--$3.3\times$ over one block at 0.95, without ensuring cross-system wins.

Null-call medians are 0.012--0.072\,ms and remain included. CAGRA uses Python, JHQ/RaBitQ C++; warm-ups differ (\cref{sec:a-t4-t5}).